\capitulo{4}{Two Ways to Read the Same Table}
Two heads on the table, the letter still hidden. For R\$\,8.80 I'll sell you a ticket that pays ten reais if the coin is A and nothing if it is B. We can do this many times. If you care about your average gain, do you buy?

The person who brought the generator reaches for a wallet. At $9/10$, the ticket is worth nine reais. The person who brought the notebook counts twenty-four A's out of twenty-eight. By that estimate the ticket is worth R\$\,8.57. The wallet stays in the pocket. If all I'd asked for was a letter, both would have said A.

Lay the training notebook open on the table.
\begin{center}
\begin{tabular}{lrrrr}
\toprule
Revealed type & 0 heads & 1 head & 2 heads & Total\\
\midrule
A & 4 & 12 & 24 & 40\\
B & 24 & 12 & 4 & 40\\
\midrule
Total & 28 & 24 & 28 & 80\\
\bottomrule
\end{tabular}
\end{center}

Put your finger on the $24$ in the first row. What else do you need to read to turn it into a probability? If you want the fraction of A rounds that gave two heads, read across the row and divide by forty. If you want the probability that the coin is A given two heads, read down the column and divide by twenty-eight. Your finger hasn't moved. The question chose the denominator.

The generator counted the other columns too and estimated each coin's chance of heads. In forty rounds of each type, it expects $40(3/4)^2=22.5$ A rounds with two heads and $40(1/4)^2=2.5$ B rounds with two heads. The column it expects to see suggests $22.5/(22.5+2.5)=9/10$. The one we saw yesterday gives $24/(24+4)=6/7$. Didn't we just insist on a single price for the same ticket? Before we toss, consider a ticket paying one real for A \emph{and} two heads. In yesterday's notebook we can count by row or by column,
\[
        \frac{40}{80}\frac{24}{40}
        =\frac{24}{80}
        =\frac{28}{80}\frac{24}{28}.
\]
In the notebook the generator expects to produce, the two prices agree as well,
\[
        \frac12\left(\frac34\right)^2
        =\frac{9}{32}
        =\frac{5}{16}\frac9{10}.
\]
The first calculation uses observed frequencies. The second builds the pair of tosses by assuming independence. To catch either learner contradicting itself, I'd have to get two prices from \emph{its own} assumptions. To see which one makes the better purchase, I need fresh rounds.

\secao{A bill that depends only on the letters}
Tomorrow, whenever I have to give an answer, I'll already have seen both tosses. I can use the training rounds to fit the probability of the letter directly for each outcome I might be shown. Learning to predict $Y$ from $X$ this way is called \emph{discriminative learning}. To see how it works, keep just the twenty-eight tickets in the last column and pick a probability $q$ to announce for A. What would the notebook charge you for that choice?

Suppose you announced $q=0.80$ for A. Make two piles of tickets, the twenty-four A's on one side and the four B's on the other. Now give A an extra hundredth of probability. In the bigger pile, the charge on each of the twenty-four tickets goes down a little. Where did that hundredth come from? Out of the $0.20$ you had set aside for B. The four tickets in that pile get more expensive. If you keep taking from B, which pile will eventually have the greater say in the bill?

Use the scale from the last chapter. For a tiny nudge in $q$, the saving on each A depends on the size of the nudge relative to the bet $q$. For a B ticket, the extra charge depends on the nudge relative to what was left, $1-q$. At $q=0.80$, the twenty-four savings pull with a strength of $24/0.80=30$. The four extra charges pull back with a strength of $4/0.20=20$. It pays to go up. Try $q=0.90$. Now the savings add up to about $26.7$ per unit of movement and the extra charges add up to 40. You've gone too far.

We can write down what we just weighed. Each A costs $-\ln q$ and each B costs $-\ln(1-q)$. Using natural logarithms just changes the units on the bill. This time the height of the curve is an expense. We're looking for the bottom of the valley. If $D$ is the total expense, its slope is
\[
 D(q)=-24\ln q-4\ln(1-q),\qquad
 D'(q)=-\frac{24}{q}+\frac{4}{1-q}.
\]
At the balance point, $24/q=4/(1-q)$. There are six times as many A tickets as B tickets. For the two sides to weigh the same, A's share of the probability must be six times B's. Six parts out of seven gives $q=6/7=24/28$. Up to that point a nudge to the right reduces the bill. Beyond it, the bill goes up. Now we know why counting the letters gives precisely the cheapest bet in this notebook.

We can do the same in the other columns. With zero heads, the cheapest bet on A is $4/28=1/7$. With one head, it is $12/24=1/2$. Our program would fit on the back of a large coin. Keep three fractions and look up the one for the column you've been given. But suppose I cover an entire column before training. You'll have no fraction to look up for that column. What would you have to add to venture an answer anyway?

\secao{How can we guess what's in the covered column?}
Let's look for a clue in the generator. With rates $3/4$ and $1/4$, each head triples the odds in favor of A. Each tail divides them by three. Take one head and one tail and cross them both off the page. Your bet is back where it started. After you've crossed off every such pair, what's left? With $H$ heads and $T$ tails, the difference $s=H-T$ counts the multiplications by three that remain, or the divisions if it's negative. Starting from odds $r_0$, we end up at $r_0\,3^s$.

Imagine the marks $1,3,9,27$ on a ruler for the odds. It costs exactly one extra head to move from one mark to the next. If I want equal steps in the notebook to cover equal distances on the ruler, where should I put those marks? Take their logarithms and they land at $0,\ln 3,2\ln 3,3\ln 3$. The difference now counts steps, each of length $\ln 3$. For a bet $q$ on A, our final position is
\[
             \ln\frac{q}{1-q}=\ln r_0+s\ln 3 .
\]
A straight line appeared before we ever mentioned regression. The starting mark came from the odds before we saw the tosses. The step length came from the difference between the two coins.

Let's keep the ruler and let the notebook choose those two measurements. Call the starting mark $\beta_0$ and the step length $\beta_1$. For a difference $s$ we read $v=\beta_0+\beta_1s$. How do we get from a mark back to a probability? If it reads $\ln 3$, there are three parts A for one part B. The bet on A is three of the four parts. If the mark reads $v$, undoing the logarithm gives $e^v$ parts A for one part B. Divide by the total number of parts. Dividing the numerator and denominator by $e^v$ just counts the same parts in a different unit. This gives us \emph{logistic regression},
\[
              q(s)=\frac{e^v}{1+e^v}=\frac{1}{1+e^{-v}},
              \qquad v=\beta_0+\beta_1s .
\]

The labels will choose the coefficients. We supplied the requirement that the three bets fit on a single straight line on this scale.\footnote{A. Y. Ng and M. I. Jordan examine what happens when we fit classifiers through the joint distribution or directly through the conditional distribution in \href{https://proceedings.neurips.cc/paper_files/paper/2001/file/7b7a53e239400a13bd6be6c91c4f6c4e-Paper.pdf}{\emph{On Discriminative vs. Generative Classifiers}} (2001).}

You've seen the numbers in the last column. Cover it anyway. We want to know whether we were still free to choose the bet in that column once we'd agreed to use the line. The two visible columns have differences of $-2$ and $0$. In the first, the odds in favor of A are four to twenty-four, or one to six. In the second they're twelve to twelve, or one to one. Moving two units to the right multiplied the odds by six. Before you lift your finger, what will you write in the last column? We're moving another two units. If the log odds lie on a straight line, we have to multiply by six again. One to one becomes six to one, giving a probability of $6/7$.

You can lift your finger. The observed odds were indeed twenty-four to four. In this table the line through the two visible points also minimizes the loss in the third column. The height is zero at $s=0$ and $\ln 6$ at $s=2$, so the fitted coefficients are $\widehat\beta_0=0$ and $\widehat\beta_1=(\ln 6)/2$. With the rates fitted by the generator, the slope was $\ln 3$. The two programs arrived at curves in the same family. One fitted rates that explain the tosses within each coin type. The other fitted the probabilities that the line assigns to the letters.

We could also use Bayes to learn the line. Whenever a letter is revealed, we multiply the weight of each candidate line by the probability it gave that letter at that value of $s$. We then normalize across candidates. The training is still discriminative. The lines bet on the letter given the tosses.

Let's do a little damage to the notebook. Correct a single letter in the middle column from B to A. The bets in the outer columns still give heights of $-\ln 6$ and $\ln 6$, but the odds in the middle are now thirteen to eleven. No straight line can meet all three requests at once. If we fit a single line to the eighty rounds, a correction in the middle can change the bet even in the two-heads column, where no letter changed. The same assumption that lets one column teach us about another can lead us astray when the relationship we imagined breaks down.

\begin{quadro}{Before you buy the ticket}
\itemq{Read the denominator} Fitting the joint distribution lets us ask how each type produces the tosses. Fitting the conditional directly lets us ask which letter to expect given the tosses we've seen. For classification, both can arrive at the same shape of curve.
\itemq{Look for what was borrowed} The generator linked tosses through a common rate for each type. Logistic regression linked columns through the slope of a line. When an observation changes the prediction for a different input, look for the assumption that linked them.
\end{quadro}

I'd seal both predictions and set aside whole rounds with fresh coins for testing. To decide whether to buy at R\$\,8.80, I want to know whether A occurs in more than $88\%$ of rounds with two heads. To assess the probabilities throughout the game, I want the average log loss on new rounds. If I decided from the notebook used for training, I'd be letting the directly fitted model grade its own exam.

So far, one head already favored A. If the first toss on its own tells us nothing and the second on its own tells us nothing either, is there any point in keeping both?

\begin{quadro}{What to take away from this chapter}
\itemq{1} Fitting the joint distribution and fitting the probability of the letter given the faces are different tasks. Both approaches can produce classifiers.
\itemq{2} With a separate, unconstrained bet for each column, logarithmic loss is minimized at the observed share of each letter.
\itemq{3} A straight line in log odds connects the columns: it lets us predict a column with no examples, but also spreads the effects of a poor assumption.
\itemq{What changed} We moved from fitting a recipe for generating rounds to fitting the answer to the question directly. We can identify the assumption that lets one example teach us something about another input.
\itemq{To carry forward} If neither face alone tells us anything about the letter, what can the two together teach us?
\end{quadro}
